\section*{अध्याय \ref{ch:g3:division} --- बाँटना और भाग}

\begin{solution}{exo:g3:division:1}
$36 \div 4 = 9$ (क्योंकि $4 \times 9 = 36$); \quad
$63 \div 9 = 7$; \quad $40 \div 5 = 8$; \quad $56 \div 8 = 7$।
\end{solution}

\begin{solution}{exo:g3:division:2}
$17 = 5 \times 3 + 2$: \omterm{def:g3:division:remainder}{भागफल} $3$, \omterm{def:g3:division:remainder}{शेषफल} $2$।

$30 = 4 \times 7 + 2$: \omterm{def:g3:division:remainder}{भागफल} $7$, \omterm{def:g3:division:remainder}{शेषफल} $2$।

$50 = 6 \times 8 + 2$: \omterm{def:g3:division:remainder}{भागफल} $8$, \omterm{def:g3:division:remainder}{शेषफल} $2$।
\end{solution}

\begin{solution}{exo:g3:division:3}
$28 = 6 \times 4 + 4$: हर बच्चे को $4$ कंचे मिलते हैं, और $4$
कंचे बच जाते हैं।
\end{solution}

\begin{solution}{exo:g3:division:4}
$68$ का \omterm{def:g3:division:half}{आधा}: $34$। $27$ की तिहाई: $9$। $48$ की \omterm{def:g3:division:half}{चौथाई}: $12$।
$150$ का \omterm{def:g3:division:half}{आधा}: $75$।
\end{solution}

\begin{solution}{exo:g3:division:5}
``$45 \div 5 = 9$'': सही, $5 \times 9 = 45$।

``$52 \div 6 = 8$, \omterm{def:g3:division:remainder}{शेषफल} $4$'': सही, $6 \times 8 + 4 = 52$ और
$4 < 6$।

``$70 \div 8 = 9$, \omterm{def:g3:division:remainder}{शेषफल} $2$'': ग़लत --- $8 \times 9 + 2 = 74$,
$70$ नहीं। सही: $70 = 8 \times 8 + 6$।
\end{solution}

\begin{solution}{exo:g3:division:6}
$32 \div 4 = 8$ पन्ने। यह \emph{समूह बनाना} है: तस्वीरें $4$-$4$
के समूहों में रखी जाती हैं, और हम समूह गिनते हैं।
\end{solution}

\begin{solution}{exo:g3:division:7}
$75 = 9 \times 8 + 3$: आठ पूरे टुकड़े, और $3$ cm का बचा हुआ
टुकड़ा।
\end{solution}

\begin{solution}{exo:g3:division:8}
$50 = 4 \times 12 + 2$: बारह कारें $48$ बच्चों को ले जाती हैं, और
$2$ बच्चों को अब भी सवारी चाहिए। $13$ कारें चाहिए।
\end{solution}

\begin{solution}{exo:g3:division:9}
$? = 7 \times 6 = 42$। \quad
$54 \div 9 = 6$, इसलिए ग़ायब भाजक $9$ है। \quad
जिसका \omterm{def:g3:division:half}{आधा} $35$ है, वह संख्या $70$ है।
\end{solution}

\begin{solution}{exo:g3:division:10}
$47 = 5 \times 9 + 2$: हर दोस्त को $9$ स्टिकर मिलते हैं, और
अमीना \omterm{def:g3:division:remainder}{शेषफल}, यानी $2$ स्टिकर रख लेती है।
\end{solution}

\begin{solution}{exo:g3:division:11}
$6$ का भाग देने पर \omterm{def:g3:division:remainder}{शेषफल} हमेशा $6$ से छोटा होता है: वह
$0$, $1$, $2$, $3$, $4$ या $5$ हो सकता है। \omterm{def:g3:division:remainder}{भागफल} $7$ वाली
सबसे बड़ी संख्या सबसे बड़ा \omterm{def:g3:division:remainder}{शेषफल} लेती है: $6 \times 7 + 5 = 47$।
\end{solution}
